\label{sec:implementation}

\subsection{Priority Queue Implementation}

The \textsc{bisect-apportion} crate implements Huntington-Hill via
a max-priority queue using Rust's \texttt{BinaryHeap}.
The priority comparison uses \emph{exact integer arithmetic} to avoid
floating-point rounding errors in tie-breaking:
\[
  P_i(s_i) = \frac{p_i}{\sqrt{s_i(s_i+1)}}
\]
is represented as the rational $p_i^2 / (s_i(s_i+1))$,
with numerator and denominator stored as \texttt{u128} integers.
Two priorities are compared by cross-multiplication:
$P_i > P_j \iff p_i^2 \cdot s_j(s_j+1) > p_j^2 \cdot s_i(s_i+1)$.
This eliminates all floating-point approximation from the priority
comparison, ensuring that tie-breaking is deterministic and exact.

The \texttt{u128} type holds values up to approximately $3.4 \times 10^{38}$.
The maximum population-squared value for the 2020 Census is
$(39{,}538{,}223)^2 \approx 1.56 \times 10^{15}$, multiplied by the
maximum seat-product $435 \times 436 \approx 1.9 \times 10^5$,
giving a maximum comparison value of approximately $3 \times 10^{20}$,
well within \texttt{u128} range.

\subsection{The 2020 Apportionment}

Running \texttt{bisect apportion --year 2020 --method huntington-hill}
produces the following seat counts (selected states):

\begin{center}
\begin{tabular}{lr|lr}
\toprule
State & Seats & State & Seats \\
\midrule
California  & 52 & Missouri       & 8 \\
Texas       & 38 & Maryland       & 8 \\
Florida     & 28 & Wisconsin      & 8 \\
New York    & 26 & Minnesota      & 8 \\
Pennsylvania& 17 & Colorado       & 8 \\
Illinois    & 17 & South Carolina & 7 \\
Ohio        & 15 & Alabama        & 7 \\
Georgia     & 14 & Kentucky       & 6 \\
North Carolina& 14 & Oregon       & 6 \\
Michigan    & 13 & Oklahoma       & 5 \\
\midrule
\multicolumn{2}{c}{\ldots} & Wyoming & 1 \\
\bottomrule
\end{tabular}
\end{center}

\noindent
Total: 435 seats.
Full table in Appendix~A.

\subsection{SHA-256 Verification}

The Census Bureau published the official 2020 apportionment on
April~26, 2021 \citep{census2020apportionment}.
The verification protocol:
\begin{enumerate}
  \item Sort the 50 (state, seats) pairs alphabetically by state name.
  \item Concatenate as ``\texttt{StateName:seats\textbackslash n}'' for each pair.
  \item Compute SHA-256 of the resulting UTF-8 string.
  \item Compare to the SHA-256 of the same string constructed from
    the Census Bureau table.
\end{enumerate}

The \texttt{bisect apportion --year 2020 --method huntington-hill
--verify} command performs this comparison automatically and exits
non-zero if the hashes differ.
The hashes match exactly, confirming zero discrepancies across
all 50 states.

\subsection{The Last Seat and the Near-Miss}

In the 2020 apportionment, the 435th seat (after the initial 50 seats
for the constitutional floor) was awarded to New York based on its
priority value at the relevant seat count.
The state that came closest to receiving an additional seat was
Montana (which received 2 seats) or New York (whose last seat was
very close to the priority cutoff).
The Census Bureau publishes the full priority sequence showing the
exact priority values at each step.

\subsection{Test Suite}

The implementation is covered by three test levels:
\begin{itemize}
  \item \textbf{L0} (unit): Priority formula produces exact values for
    small inputs; $P(1) = p/\sqrt{2}$, $P(2) = p/\sqrt{6}$, etc.
  \item \textbf{L1} (integration): 5-state synthetic examples
    produce expected seat counts that agree with hand computation.
  \item \textbf{L2} (regression): 2020, 2010, and 2000 Census data
    produce seat counts matching official Census Bureau apportionments
    at 100\% match rate.
\end{itemize}
